A. FHM model with different pulse durations 

==================
B. Satellites
- 2 satellites in 2 orbits. Each has 2 variables (see FM paper)
one of the satellites is on a fixed orbit.
the other satellite, transitions from orbit 1 -> orbit 2 -> orbit 3 based on some timing. 
- Safety property show that the 2 satellites are close at the end.

The above 2 are easy to create but lower impact 

==================

C. Cars

We will create models of cars moving on a straight highway or taking turns through intersections. The models will be modular. That is, by putting modules for multiple cars, we will be able to create more complicated models (easily).  

One thing to keep in mind: If we lose the sense of global time in the verification, then does it still make sense to talk about properties that measure distance between two cars $|x(t) - y(t)|$? 

Models. Keep it simple for now, e.g., Dubin's vehicles. Each car automata has the following variables:

x,y: positions
v: forward velocity
\theta: heading w.r.t positive x axis
\omega: speed of turning

Each car has the following parameters:

l: length, w: width
a1,a2,a0: forward acceleration/braking 
aw: turn rate control

Each car has the following modes:

S0. straight (no turn, no speed-up could be along and direction)
\dot x = v cos(\theta)
\dot y = v sin(\theta)
\dot v = 0
\dot \theta = 0


S1,2. straight speed-up (no turn,  speed-up fast or slow)
same as S0 except
\dot v = a1 or a2

should not be in this mode for too long; upper speed bound

S3. straight slowdown (no turn,  slow-down )
same as S0 except
\dot v = a3 < 0

should not be in this mode for too long; lower speed bound

T1,2. simple left and right turns (no speed-up)
\dot x = v cos(\theta)
\dot y = v sin(\theta)
\dot v = 0
\dot \theta = +/-aw

By doing timed transition across each of these modes we can create 1-car automata modeling following maneuvers

1. straight drive S1/S2/S3
2. lane shift S1->T1->T2->S1
3. left turn S1->T1->S1
4. ...

One we create these modules (nicely), we can then go to town and develop a whole array of interesting composed automata that capture interesting verification scenarios. Combine a couple of cars in a model to create scenario automata. Each of these scenarios can be made safe/unsafe by tweaking 1/2 parameters. Here are a few examples:

1. Car A going straight (constant speed mode S1) with some initial position and velocity range for x1, v1. Then slows down at time t1 (+/- some range) then going constant speed S1 again. Car B starts distance at some x2 behind x1 with some velocity range. Goes constant speed for a while and then starts to brake at time t1 + t2. What is the max t2 that keeps B safe?

2. Same scenario as above, except: Car B tries to switch to a lane at time t1 + t2, to avoid car A.

3. Car A going straight (constant speed S1) with some initial position and velocity range for x1, v1. Car B overtakes car A based on timed switches S1->T1->T2->S2->T2->T1->S1 (like the model we have already). What is the smalled duration of time Car B spends in left-lane?

4. Changing lane between two cars 

5. Taking a right and merging between two cars

6. Taking a left and merging between two cars